\section{Promise dan async/await}

\textbf{Promise} merepresentasikan nilai yang mungkin tersedia di masa depan. Promise memiliki tiga state: \textit{pending} (menunggu), \textit{fulfilled} (berhasil, dengan nilai), dan \textit{rejected} (gagal, dengan error). \code{.then()} menangani fulfillment; \code{.catch()} menangani rejection; \code{.finally()} dijalankan di kedua kasus. Promise dapat di-chain untuk operasi berurutan \cite{jsInfo}.

\begin{table}[htbp]
\centering
\begin{tabular}{|P{3cm}|P{4.5cm}|P{4.5cm}|}
\hline
\textbf{Pola} & \textbf{Sintaks} & \textbf{Kelebihan} \\
\hline
Callback & \code{fn(data, callback)} & Sederhana; tapi nesting \\
\hline
Promise & \code{fn().then().catch()} & Chain-able; error propagation \\
\hline
async/await & \code{const r = await fn()} & Linier seperti sinkron \\
\hline
\end{tabular}
\caption{Evolusi penanganan asinkron di JavaScript}
\end{table}

\begin{webcode}[language=JavaScript, caption={Promise dan async/await}]
// Promise
function fetchUser(id) {
  return fetch(`/api/users/${id}`)
    .then(res => res.json());
}
fetchUser(1).then(user => console.log(user.name));

// async/await (cara modern, lebih mudah dibaca)
async function getUser(id) {
  try {
    const res = await fetch(`/api/users/${id}`);
    const user = await res.json();
    console.log(user.name);
  } catch (err) {
    console.error("Gagal:", err.message);
  }
}
\end{webcode}

\begin{konsep}
\textbf{async/await = Sugar untuk Promise.} \code{async function} selalu mengembalikan Promise. \code{await} menunggu Promise selesai dan mengembalikan nilainya. Error ditangkap dengan \code{try/catch} biasa (bukan \code{.catch()}). Ini membuat kode asinkron terlihat dan berperilaku seperti kode sinkron---jauh lebih mudah dibaca dan di-debug dibanding \code{.then()} chain \cite{mdnJsGuide}.
\end{konsep}

